% Paper 4, §5 — The correlation–causation gap, and how the field closed it
% Thesis: high-accuracy probes are *correlational*. Hewitt & Liang 2019
% surfaced the methodology critique with control tasks; Belinkov 2022
% consolidated it into a taxonomy of confounders (probe expressivity,
% spurious correlation, vestigial information). The 2023+ resolution:
% pair probes with activation interventions — direction-ablation tests
% causal use, direction-addition (the activation-steering family) turns
% probe directions into behavioral controls. Marks–Tegmark 2023 is the
% canonical worked example. The "patch-then-probe" pipeline composes
% §9's patching machinery with §4's probes.

\section{The correlation--causation gap}
\label{sec:correlation-causation}

A high-accuracy \glsterm{probe-2}{probe} at layer $\ell$ tells us that property $y$ is
\emph{linearly decodable} from $\mathbf{h}^{(\ell, t)}$. It does not
tell us that the model \emph{uses} the decoded direction to compute its
output \citep{belinkov2022-probing-survey}. The 2017--2019 wave of
probing methodology summarized in \S\ref{sec:probing} put the
correlational claim on firm footing; the 2019--2022 wave, anchored by
\citet{hewitt2019-control-tasks} and consolidated by
\citet{belinkov2022-probing-survey}, made the gap between
\emph{decodable} and \emph{used} into a methodological problem the
field had to confront. The 2023+ wave closes the gap by pairing
probes with two activation-level interventions: \emph{direction
ablation} along the \glsterm{probe-2}{probe} direction $\hat{\mathbf{d}}_y$, which tests
whether removing $\mathbf{d}_y$ from the \glsterm{residual-stream}{residual stream} changes
behavior, and \emph{direction addition} $\mathbf{h} \mapsto \mathbf{h} +
\alpha \hat{\mathbf{d}}_y$, the activation-steering family that turns
a \glsterm{probe-2}{probe} direction into an inference-time control surface. This section
states the critique formally, gives the two intervention operators in
canonical form, walks through the Marks--Tegmark 2023 truth-direction
worked example, and shows how the resulting ``patch-then-\glsterm{probe-2}{probe}''
pipeline composes patching (\S\ref{sec:patching}) with probing
(\S\ref{sec:probing}) into a single causal recipe.

\subsection{The Hewitt--Liang control-task critique}
\label{sec:control-tasks}

\citet{hewitt2019-control-tasks} surface the load-bearing methodology
problem: a sufficiently expressive \glsterm{probe-2}{probe} can achieve high accuracy
even when the underlying model does not represent the property in any
useful way, because the \glsterm{probe-2}{probe} itself learns the property from
low-information features in $\mathbf{h}$. The proposed fix is the
\emph{control task}: an alternative probing target with the same
input--output type as the real task but with the labels permuted to be
informationally meaningless. A part-of-speech control task, for
instance, assigns each word type a single random label drawn once and
held fixed across the dataset; the \glsterm{probe-2}{probe} is then asked to learn this
nonsense mapping. The diagnostic is the \emph{selectivity}, defined
as the gap
\begin{equation}
  \mathrm{selectivity}(g, \ell)
  \;:=\; \mathrm{acc}(g, \ell, \text{real task})
  \;-\; \mathrm{acc}(g, \ell, \text{control task}).
  \label{eq:selectivity}
\end{equation}
\citep{hewitt2019-control-tasks}. A \glsterm{probe-2}{probe} that scores high on both
real and control tasks has low selectivity: its accuracy is a property
of the \glsterm{probe-2}{probe}'s expressivity, not of the model's representation. A
\glsterm{probe-2}{probe} that scores high on the real task but low on the control task
has high selectivity, and the real-task accuracy can be attributed to
information the model itself encoded in $\mathbf{h}^{(\ell, t)}$.
\citet{hewitt2019-control-tasks} argue this should be a non-negotiable
companion measurement to any probing accuracy reported in the
literature, and Eq.~\eqref{eq:selectivity} is now the standard
control alongside the held-out test split of the probing protocol of
\S\ref{sec:probe-protocol}.

The control-task critique is sharper than ``be careful with deep
probes.'' Even a \glsterm{linear-probe}{linear probe} can fail selectivity on a poorly
initialized model, on noisy activations, or on a label set whose
distribution is dominated by frequency effects rather than the
property of interest. Conversely, a high-selectivity \glsterm{probe-2}{probe} is not yet
proof of \emph{causal} use, only of \emph{informational} encoding;
selectivity rules out \emph{\glsterm{probe-2}{probe}-expressivity} confounds, but two
further confounds remain.

\subsection{Belinkov's three confounders}
\label{sec:belinkov-taxonomy}

\citet{belinkov2022-probing-survey} consolidates the post-control-task
critique into a three-confounder taxonomy that has become the
methodological \glsterm{vocabulary}{vocabulary} of the
field\footnote{KB excerpt: \texttt{kb/excerpts/belinkov2022-probing-survey.md\#abstract}.}.
The first is \citet{hewitt2019-control-tasks}'s
\textbf{\glsterm{probe-2}{probe}-expressivity} confound, addressed by Eq.~\eqref{eq:selectivity}.
The second is the \textbf{spurious-correlate} confound: even a
selective \glsterm{probe-2}{probe} may be reading a \glsterm{feature}{feature} $z$ that correlates with
$y$ in the probing dataset, not $y$ itself. The \glsterm{probe-2}{probe} direction
$\hat{\mathbf{d}}_y$ is then a direction for $z$, and any downstream
intervention along $\hat{\mathbf{d}}_y$ is an intervention on $z$
masquerading as one on $y$. The standard mitigation is
\emph{cross-distribution generalization}: train the \glsterm{probe-2}{probe} on dataset
$\mathcal{D}_1$, evaluate on dataset $\mathcal{D}_2$ chosen to break
the $z$--$y$ correlation. \citet[\S 1, \S 4]{marks-tegmark-2023-truth}
make exactly this move the cornerstone of their truth-direction
analysis, training on \texttt{cities} and evaluating on
\texttt{neg\_cities}, \texttt{cities\_cities\_conj}, and the
\texttt{likely} control of nonfactual completions designed to match
text probability across the true/false split.

The third is \textbf{vestigial information}: the model may have
computed $y$ at some intermediate stage, but later layers may not use
the resulting direction; the \glsterm{probe-2}{probe} reads the trace, not the driver.
This is the most insidious confound because it is invisible to both
selectivity and cross-distribution generalization --- a \glsterm{probe-2}{probe} of a
truly vestigial \glsterm{feature}{feature} can be high-selectivity \emph{and} robust
across distributions, because the \glsterm{feature}{feature} really is encoded by the
model on those distributions. The only way to distinguish a vestigial
encoding from a used one is intervention: change the activation,
hold the rest of the model fixed, observe whether the output behavior
moves. \citet{belinkov2022-probing-survey}'s ``advances'' section
catalogs the three intervention families that close this gap ---
amnesic probing \citep[Elazar et al.\ 2021, surveyed in
][]{belinkov2022-probing-survey}, causal-mediation analysis (the
ancestor of the activation-patching protocol of
\S\ref{sec:patching}), and interchange interventions --- and the
2023+ probing literature is the inheritor of all three.

\begin{contradiction}
\textbf{Decodable does not entail used.} A \glsterm{probe-2}{probe} with selectivity
gap of $\Delta = 0.9$ on a held-out test split, generalizing across
five datasets that break the obvious surface confounders, is still
\emph{correlational}. The model may compute the property at some
intermediate layer and discard it; later layers may use a different
representation; the \glsterm{probe-2}{probe} direction may correlate near-perfectly with
the used representation while not being it. The strongest form of the
critique \citep[\S abstract]{belinkov2022-probing-survey}: probing
methodology has historically conflated ``$y$ is in $\mathbf{h}$''
with ``the model uses $y$ in $\mathbf{h}$,'' and the field's record
of behavioral interventions that did \emph{not} work after a
\glsterm{probe-2}{probe}-only validation is precisely what motivated the 2023+ pivot to
intervention-paired probing in \S\ref{sec:ablate-add}.
\end{contradiction}

\subsection{Direction ablation and direction addition}
\label{sec:ablate-add}

The 2023+ resolution is to compose probing with one of two
activation-level interventions, both of which act linearly on the
\glsterm{residual-stream}{residual stream} along a unit-norm \glsterm{probe-2}{probe} direction
$\hat{\mathbf{d}}_y \in \mathbb{R}^D$.

\paragraph{Direction ablation.} Project the residual-stream activation
onto the orthogonal complement of $\hat{\mathbf{d}}_y$:
\begin{equation}
  \boxed{\;
  \mathbf{h} \;\mapsto\;
  \mathbf{h} \;-\; (\mathbf{h}^\top \hat{\mathbf{d}}_y)\,\hat{\mathbf{d}}_y
  \;=\; \big(\mathbf{I} - \hat{\mathbf{d}}_y \hat{\mathbf{d}}_y^\top\big)\,\mathbf{h}.
  \;}
  \label{eq:direction-ablation}
\end{equation}
The operator on the right is the orthogonal projector
$\mathbf{P}_{\hat{\mathbf{d}}_y^{\perp}} :=
\mathbf{I} - \hat{\mathbf{d}}_y \hat{\mathbf{d}}_y^{\top}$ onto the
$(D-1)$-dimensional subspace orthogonal to $\hat{\mathbf{d}}_y$
\citep{marks-tegmark-2023-truth}. Inserting Eq.~\eqref{eq:direction-ablation}
into the \glsterm{residual-stream}{residual stream} at one or more layers and re-running the rest
of the model removes whatever signal $\hat{\mathbf{d}}_y$ carries,
leaving every other direction unchanged in the activation. The
diagnostic question is then simple: if the model's output behavior on
property $y$ changes substantially under ablation, $\hat{\mathbf{d}}_y$
is causally implicated; if behavior is unchanged, the direction is
vestigial or correlated with a separate driver. Eq.~\eqref{eq:direction-ablation}
is the formal closing move on the vestigial-information confound of
\S\ref{sec:belinkov-taxonomy}: it does not just check whether
information is encoded along $\hat{\mathbf{d}}_y$, it checks whether
the model's downstream computation \emph{depends} on it.

The shape ledger: at one batched residual-stream slice the activation
is $\mathbf{h} \in \R^{\BSH}$, the
direction is $\hat{\mathbf{d}}_y \in \mathbb{R}^D$, the projection
$\mathbf{h}^\top \hat{\mathbf{d}}_y$ contracts the $D$-axis to a
scalar per $(B, S)$ position, and the ablation rewrites $\mathbf{h}$
in place. The intervention is one matrix subtraction per layer where
it is applied; cost is dominated by re-running the suffix
$\mathcal{M}_{>\ell}$ from the modified state.

\paragraph{Direction addition (activation steering).} Add a scaled
copy of $\hat{\mathbf{d}}_y$ to the \glsterm{residual-stream}{residual stream}:
\begin{equation}
  \boxed{\;
  \mathbf{h} \;\mapsto\;
  \mathbf{h} \;+\; \alpha\,\hat{\mathbf{d}}_y,
  \;}
  \label{eq:direction-addition}
\end{equation}
with $\alpha \in \mathbb{R}$ a steering coefficient that can be
positive (push the model toward $y = +$) or negative (push toward
$y = -$) \citep{marks-tegmark-2023-truth}. Direction addition is the
method underneath the activation-steering family of
\deepencite{zou2023-representation-engineering §3, citation-pending}
(representation engineering, which casts steering as a control-theoretic
intervention on the \glsterm{residual-stream}{residual stream}'s representational subspace) and
\deepencite{rimsky2024-caa §3, citation-pending} (Contrastive
Activation Addition, which derives $\hat{\mathbf{d}}_y$ from a
difference-of-means over contrastive prompt pairs and applies it at a
chosen layer). The 2025 generalization to behaviorally-singular
properties that decompose into multi-axis subspaces appears in
\deepencite{vennemeyer2025-sycophancy-not-one-thing §3,
citation-pending}, who report that \glsterm{sycophancy}{sycophancy} in current chat models
factorizes into three independently steerable directions (agreement,
praise, genuine-agreement), each detected by \glsterm{mass-mean-probing}{mass-mean probing} and
controllable by Eq.~\eqref{eq:direction-addition} at the
single-direction level.

The relationship between the two operators: ablation in
Eq.~\eqref{eq:direction-ablation} is the limit of subtraction along
$\hat{\mathbf{d}}_y$ that takes the projection coefficient to zero,
while addition in Eq.~\eqref{eq:direction-addition} is a free shift
along $\hat{\mathbf{d}}_y$ leaving the projection of the original
$\mathbf{h}$ in place. Ablation tests \emph{necessity} (does removing
the direction break behavior?); addition tests \emph{sufficiency}
(does adding the direction induce behavior?). A causally-implicated
direction passes both tests; a spurious or vestigial direction passes
neither, or passes only one in a way that fails to compose.

\begin{intuition}
The ablation operator
$\mathbf{P}_{\hat{\mathbf{d}}_y^{\perp}} = \mathbf{I} -
\hat{\mathbf{d}}_y \hat{\mathbf{d}}_y^{\top}$ in
Eq.~\eqref{eq:direction-ablation} is the same orthogonal projector
that appears in Gram--Schmidt orthogonalization: subtracting from
$\mathbf{h}$ its component along $\hat{\mathbf{d}}_y$ leaves a vector
in the hyperplane normal to $\hat{\mathbf{d}}_y$. Causally probing
along $\hat{\mathbf{d}}_y$ is, in this sense, asking whether the
model's downstream computation lives \emph{entirely in that
hyperplane}: if so, the projection is a no-op and behavior is
unchanged; if the computation steps out of the hyperplane along
$\hat{\mathbf{d}}_y$ to produce its answer, the projection blocks
that step. Returning to canonical form: Eq.~\eqref{eq:direction-ablation}
is rank-$(D-1)$, the lost rank is exactly $\mathrm{span}(\hat{\mathbf{d}}_y)$,
and the experiment is whether the model's output changes when this
rank-1 subspace is zeroed out at the \glsterm{residual-stream}{residual stream}.
\end{intuition}

\subsection{The Marks--Tegmark worked example: flipping TRUE/FALSE}
\label{sec:mt-worked-example}

\citet{marks-tegmark-2023-truth} provide the canonical worked example
of the resolution. The mass-mean direction
$\mathbf{d}_{\mathrm{MM}} = \boldsymbol{\mu}_+ - \boldsymbol{\mu}_-$
introduced in \S\ref{sec:mass-mean} (Eq.~4 of \S\ref{sec:probing}) is
fit on the \texttt{cities} dataset of LLaMA-2-13B residual-stream
activations at the period-token of layer 12. \glsterm{probe-2}{Probe} accuracy on the
held-out test split is high; cross-dataset transfer to
\texttt{neg\_cities}, \texttt{larger\_than}, and \texttt{cities\_cities\_conj}
is sufficient to argue that $\mathbf{d}_{\mathrm{MM}}$ is not a
surface artifact of the prompt template
\citep[\S 1.1, \S 4]{marks-tegmark-2023-truth}\footnote{KB excerpt:
\texttt{kb/excerpts/marks-tegmark-2023-truth.md\#sec-1-related},
\texttt{\#sec-4-misalignment}.}. The closing causal step is the
ablation experiment from Eq.~\eqref{eq:direction-ablation}: the
authors apply
$\mathbf{P}_{\hat{\mathbf{d}}_{\mathrm{MM}}^\perp}$ at the localized
hidden states identified by the patching pipeline of
\S\ref{sec:patching} and report:
\begin{quote}
Causal evidence obtained by surgically intervening in a \glsterm{llm}{LLM}'s forward
pass, causing it to treat false statements as true and vice
versa.\hfill\citep[abstract]{marks-tegmark-2023-truth}\footnote{KB
excerpt: \texttt{kb/excerpts/marks-tegmark-2023-truth.md\#abstract}.}
\end{quote}
The TRUE/FALSE flip is the load-bearing observation. The \glsterm{probe-2}{probe} alone
identifies a direction that linearly separates the two classes
(\S\ref{sec:probe-setup}); cross-dataset generalization rules out the
obvious spurious-correlate failure modes (\S\ref{sec:belinkov-taxonomy});
the ablation in Eq.~\eqref{eq:direction-ablation} demonstrates that
removing $\mathbf{d}_{\mathrm{MM}}$ from the localized hidden states
causes the model to systematically misclassify factual statements
that the unmodified model gets right. By the necessity criterion of
\S\ref{sec:ablate-add}, the mass-mean direction is causally
implicated; the model is not just \emph{representing} truth, it is
\emph{using} the representation. The same paper reports that
direction-addition (Eq.~\eqref{eq:direction-addition}) along
$\mathbf{d}_{\mathrm{MM}}$ steers the model's outputs in the predicted
direction, providing the matching sufficiency
test\footnote{KB excerpt: \texttt{kb/excerpts/marks-tegmark-2023-truth.md\#sec-1-contributions}.}.

The strongest form of the conclusion --- ``$\mathbf{d}_{\mathrm{MM}}$
is the \glsterm{truth-direction-2}{truth direction}'' --- is, however, qualified by the misalignment
finding of \S\ref{sec:truth-directions} (forward link): even the
mass-mean direction is dataset-dependent in early-mid layers, with
\texttt{cities} and \texttt{neg\_cities} truth directions
\emph{antipodal} in early LLaMA-2-13B layers and only aligned by
late layers \citep[\S 4]{marks-tegmark-2023-truth}\footnote{KB excerpt:
\texttt{kb/excerpts/marks-tegmark-2023-truth.md\#sec-4-misalignment}.}.
The causal-validation argument lifts the load-bearing claim from
``decodable'' to ``used'' but does not, by itself, resolve whether
``truth'' is a single linear axis or a multi-axis subspace; the next
section takes up that geometry directly.

\subsection{The patch-then-probe pipeline}
\label{sec:patch-then-probe}

The Marks--Tegmark causal pipeline is more than a \glsterm{probe-2}{probe} with an
ablation appended; it is a two-stage composition that uses the
patching machinery of \S\ref{sec:patching} to choose \emph{where} to
\glsterm{probe-2}{probe}. The first stage is \glsterm{activation-patching}{activation patching}: sweep the
(component, token, layer) grid of $\mathrm{AIE}_C$ values
(Eq.~\ref{eq:aie} of \S\ref{sec:patching-protocol}) on a
true/false counterfactual benchmark, identify the small group of
hidden states with high causal mediation, and restrict the probing
domain to those states. \citet[\S 3]{marks-tegmark-2023-truth} report
three groups of causally-implicated hidden states for LLaMA-2-13B on
the \texttt{cities} dataset; the third (group \emph{c}) reads
directly as TRUE/FALSE under the model's own decoder head, while the
second (group \emph{b}) sits over the final token of the statement
and end-of-sentence punctuation and ``store[s] a representation of
the statement's truth''\footnote{KB excerpt:
\texttt{kb/excerpts/marks-tegmark-2023-truth.md\#sec-3-layers}.}.

The second stage is probing within those states. Mass-mean
(Eq.~\ref{eq:mass-mean} of \S\ref{sec:probing}) is fit on the
group-\emph{b} activations rather than on the full grid, and the
ablation of Eq.~\eqref{eq:direction-ablation} is applied at the same
locations. The composition is what makes the causal argument
load-bearing: patching localizes the components that carry the
behavior, probing localizes the direction \emph{within} those
components, and ablation tests whether that direction is the one the
model uses. Without the patching stage, the probing direction is
fit on activations that may or may not be causally relevant; with it,
the \glsterm{probe-2}{probe}'s domain is restricted to states that patching has already
shown to mediate the behavior.

\begin{analogy}
Patching localizes the \emph{room}; probing localizes the \emph{shelf
within the room}; ablation tests whether the book on that shelf is
the one being read aloud. A \glsterm{probe-2}{probe} alone tells us a book about $y$ is
on \emph{some} shelf in the library; patching alone tells us a book
about $y$ is in \emph{this} room; the composition tells us \emph{this
book on this shelf in this room} is the one whose absence makes the
narrator stop reading. Returning to canonical form: patching's $C$ is
the room (a component subset of the residual-stream grid), the
\glsterm{probe-2}{probe}'s $\hat{\mathbf{d}}_y$ is the shelf (a one-dimensional
subspace of $\mathbb{R}^D$ at that component), and the ablation
$\mathbf{I} - \hat{\mathbf{d}}_y \hat{\mathbf{d}}_y^\top$ in
Eq.~\eqref{eq:direction-ablation} tests the causal entailment by
zeroing the shelf and re-running the suffix $\mathcal{M}_{>\ell}$.
\end{analogy}

The patch-then-\glsterm{probe-2}{probe} pipeline is now the de facto causal-probing
recipe in the truth-direction literature
\citep{marks-tegmark-2023-truth}, and the dual-method recipe behind
many of the cross-method convergence cases catalogued in
\S\ref{sec:cross-method}. It also exposes the dependencies between
methods: patching's STR-vs-GN corruption sensitivity
(\S\ref{sec:patching-corruption}, especially the
\texttt{[CONTRADICTION]} block on GN mis-localization) and the
metric-choice pitfalls (\S\ref{sec:patching-metric}) propagate into
the \glsterm{probe-2}{probe} domain, since the localized hidden states are the
\emph{output} of a patching pass whose recommendations matter.

\subsection{Cross-paper threads and forward link}
\label{sec:correlation-causation-forward}

Two cross-paper threads sit in the load-bearing path of this section.
The \glsterm{chain-of-thought}{chain-of-thought}-faithfulness question of Paper~3 \S 11 is, in
the framing of this section, a probing question: a behavioral \glsterm{probe-2}{probe}
of ``the model is using its written reasoning'' is correlational
unless paired with an intervention --- delete or rewrite the
intermediate reasoning trace and observe whether the final answer
changes --- in exactly the necessity-and-sufficiency form of
\S\ref{sec:ablate-add}. The \glsterm{probe-2}{probe}-on-activations and the
\glsterm{probe-2}{probe}-on-\glsterm{chain-of-thought}{CoT}-text are different observation modalities, but the
methodological pattern is the same one this section formalizes, and
the same set of confounders (\glsterm{probe-2}{probe} expressivity, spurious correlate,
vestigial information) translates across modalities. The alignment
chapter of Paper~5 consumes behavioral probes as the detection
methodology for deception, refusal-circumvention, \glsterm{sycophancy}{sycophancy}, and
\glsterm{scheming}{scheming}; every claim there of the form ``the model is being
$y$-deceptive'' depends on a \glsterm{probe-2}{probe} that has passed both the
selectivity and the ablation tests of this section, or it inherits
the correlational caveat in full.

The forward links within this paper are immediate.
\S\ref{sec:truth-directions} picks up the truth-direction case study
at the geometry-of-truth depth, asking when the same model's
mass-mean direction is dataset-stable and when it is not, and what
the misalignment in early-mid layers tells us about whether
``truth'' is a single linear axis or a multi-axis subspace.
\S\ref{sec:cross-method} returns to the patch-then-\glsterm{probe-2}{probe} pipeline as
one of the field's clearest cross-method convergence stories, with
\S\ref{sec:patching}'s machinery on the activation-patching side
and \S\ref{sec:probing}'s machinery on the probing side composing
into a causal pipeline that neither method supports alone. The
correlation--causation gap is, in this paper's framing, the cleanest
example of why method-pluralism is principled: the gap is closed not
by a better single method, but by composing two methods whose
individual commitments are insufficient and whose joint commitments
are not.
